\documentclass[11pt,a4paper]{article}

% Packages
\usepackage[margin=2.5cm]{geometry}
\usepackage[T1]{fontenc}
\usepackage[utf8]{inputenc}
\usepackage{lmodern}
\usepackage{microtype}
\usepackage{booktabs}
\usepackage{longtable}
\usepackage{hyperref}
\usepackage{graphicx}
\usepackage{xcolor}
\usepackage{listings}
\usepackage{enumitem}
\usepackage{titlesec}
\usepackage{fancyhdr}
\usepackage{amsmath}

% Hyperref setup
\hypersetup{
    colorlinks=true,
    linkcolor=blue!60!black,
    urlcolor=blue!60!black,
    citecolor=blue!60!black
}

% Header / Footer
\pagestyle{fancy}
\fancyhf{}
\fancyhead[L]{Flight Software Simulation}
\fancyhead[R]{Design Document}
\fancyfoot[C]{\thepage}
\renewcommand{\headrulewidth}{0.4pt}

% Title
\title{%
    \vspace{-1cm}
    \textbf{Flight Software Simulation}\\[0.4em]
    \Large Software-Bus Architecture for a Simplified Avionics System\\[0.6em]
    \large Design Document
}
\author{Pablo Buceta Santos}
\date{\today}

\begin{document}
\maketitle
\tableofcontents
\newpage

%----------------------------------------------------------------------
\section{Introduction \& Motivation}
%----------------------------------------------------------------------

This document describes the design of a \textbf{software-bus-based flight software simulation} developed as a personal project in preparation for a Flight Software Engineer interview at Rocket Lab (Electron).

The simulation implements a classic spacecraft avionics architecture centred on a Software Bus. Its purpose is to explore and demonstrate the interactions between typical flight-software components (command handling, telemetry, GNC, health monitoring, scheduling, etc.) in a clean, well-structured C codebase.

\textbf{Important scope statement:}
This is an \emph{architectural and functional simulation}, not a flight-ready or hard-real-time system. It prioritises clarity of component boundaries, message flows and design decisions over real-time guarantees, redundancy or certification artefacts.

%----------------------------------------------------------------------
\section{Overall Architecture}
%----------------------------------------------------------------------

The system follows a classic \textbf{software-bus} (publish/subscribe) architecture commonly found in spacecraft flight software.

\begin{center}
% Placeholder for architecture diagram
\textit{[Architecture diagram -- to be inserted]}
\end{center}

All major components communicate exclusively through a central \textbf{Software Bus}. This decouples producers from consumers and makes the data flow explicit and inspectable.

\subsection{High-level data paths}

\begin{itemize}[leftmargin=*]
    \item \textbf{Command path:} External world $\rightarrow$ Communication Interface $\rightarrow$ Command Manager $\rightarrow$ Software Bus $\rightarrow$ target component(s)
    \item \textbf{Telemetry path:} Components / Instruments $\rightarrow$ Telemetry Manager $\rightarrow$ Software Bus $\rightarrow$ Communication Interface $\rightarrow$ External world
    \item \textbf{Health path:} Components publish heartbeats / status $\rightarrow$ Health Manager $\rightarrow$ Software Bus (alerts / safe-mode flags)
    \item \textbf{Control path:} GNC subscribes to sensor data and commands, publishes control outputs and vehicle state
\end{itemize}

%----------------------------------------------------------------------
\section{Frozen Module List and Responsibilities}
%----------------------------------------------------------------------

The following module list is considered \textbf{frozen} for the first complete version of the simulation. Responsibilities may be refined in wording, but the set of modules and their primary duties will not change.

\subsection{Software Bus}
\begin{description}[style=nextline, leftmargin=1.5em]
    \item[Primary role] Central communication backbone of the system.
    \item[Responsibilities]
    \begin{itemize}
        \item Provide a publish/subscribe (or typed message) interface.
        \item Route messages from publishers to all current subscribers of a given topic / message type.
        \item Guarantee in-order delivery within a single topic for the simulation (best-effort is acceptable).
        \item Allow components to register and deregister subscriptions at start-up (dynamic subscription is optional).
        \item Provide a simple, deterministic API usable from pure C.
    \end{itemize}
    \item[Non-responsibilities] Message persistence across restarts, network distribution, prioritisation beyond simple ordering.
\end{description}

\subsection{Scheduler}
\subsubsection*{Primary role}
Drive the temporal behaviour of the simulation.
\subsubsection*{Responsibilities}
\begin{itemize}
\item Maintain a simple cyclic executive or priority-based dispatch loop.
\item Call each component's periodic ``tick'' function at a defined rate (or deliver pending messages).
\item Support both periodic tasks and aperiodic / event-driven work.
\item Provide a global simulation time source.
\item Allow the system to be stepped (useful for deterministic demos and testing).
\end{itemize}
\subsubsection*{Non-responsibilities}
Hard real-time guarantees, multi-core scheduling, deadline miss recovery.

\subsection{Command Manager}
\begin{description}[style=nextline, leftmargin=1.5em]
    \item[Primary role] Entry point and validator for all telecommands.
    \item[Responsibilities]
    \begin{itemize}
        \item Receive raw commands from the Communication Interface.
        \item Perform basic validation (opcode range, parameter length, simple sanity checks).
        \item Translate validated commands into internal bus messages.
        \item Publish command messages onto the Software Bus for the intended destination component(s).
        \item Optionally keep a short command history / log for debugging.
    \end{itemize}
    \item[Non-responsibilities] Complex authentication, encryption, or long-term command storage.
\end{description}

\subsection{Telemetry Manager}
\begin{description}[style=nextline, leftmargin=1.5em]
    \item[Primary role] Collection, packaging and egress of telemetry.
    \item[Responsibilities]
    \begin{itemize}
        \item Subscribe to telemetry-related topics on the Software Bus.
        \item Aggregate or sample data from multiple producers.
        \item Package telemetry into well-defined packets (or frames).
        \item Forward packaged telemetry to the Communication Interface.
        \item Support basic filtering or rate limiting if needed for demos.
    \end{itemize}
    \item[Non-responsibilities] On-board storage of large telemetry archives, complex CCSDS framing (can be added later).
\end{description}

\subsection{GNC (Guidance, Navigation \& Control)}
\begin{description}[style=nextline, leftmargin=1.5em]
    \item[Primary role] Maintain a simplified vehicle state and produce control outputs.
    \item[Responsibilities]
    \begin{itemize}
        \item Subscribe to relevant sensor data and telecommands.
        \item Maintain a minimal internal state vector (position, velocity, attitude, or a reduced subset).
        \item Run a simple control law or guidance algorithm (even a trivial proportional controller is acceptable for v1).
        \item Publish vehicle state and control commands onto the bus.
        \item Expose key state variables as telemetry.
    \end{itemize}
    \item[Non-responsibilities] High-fidelity orbital dynamics, Kalman filtering, or production-grade GNC algorithms.
\end{description}

\subsection{Health Manager}
\subsubsection*{Primary role}
System health monitoring and basic fault awareness.
\subsubsection*{Responsibilities}
\begin{itemize}
\item Subscribe to heartbeat / status messages from other components.
\item Detect missing heartbeats or explicit error flags.
\item Maintain an overall system health state (e.g. nominal / degraded / safe).
\item Publish health status and alerts onto the Software Bus.
\item Optionally trigger a simple safe-mode reaction (e.g. inhibit actuators).
\end{itemize}
\subsubsection*{Non-responsibilities}
Full FDIR (Fault Detection, Isolation and Recovery), redundant voting, or permanent failure logging.

\subsection{Memory Manager}
\begin{description}[style=nextline, leftmargin=1.5em]
    \item[Primary role] Controlled memory allocation and usage tracking for the simulation.
    \item[Responsibilities]
    \begin{itemize}
        \item Provide a simple deterministic allocator or memory pools for components that need dynamic memory.
        \item Track total and per-component memory usage.
        \item Publish memory statistics as telemetry.
        \item Detect and report allocation failures to the Health Manager.
    \end{itemize}
    \item[Non-responsibilities] Full heap implementation with fragmentation handling of production quality, virtual memory, or MMU simulation.
\end{description}

\subsection{Instrument Manager}
\begin{description}[style=nextline, leftmargin=1.5em]
    \item[Primary role] High-level abstraction and coordination of instruments.
    \item[Responsibilities]
    \begin{itemize}
        \item Translate bus messages into driver-level requests.
        \item Aggregate data from multiple instruments when needed.
        \item Enforce simple instrument mode management (on/off, standby, etc.).
        \item Publish instrument data and status onto the Software Bus.
    \end{itemize}
    \item[Non-responsibilities] Direct hardware register access (this belongs to the Drivers Layer).
\end{description}

\subsection{Drivers Layer}
\begin{description}[style=nextline, leftmargin=1.5em]
    \item[Primary role] Low-level simulated device drivers.
    \item[Responsibilities]
    \begin{itemize}
        \item Implement a thin interface that looks like a hardware driver (init, read, write, ioctl-style).
        \item Interact with the simulated Instruments.
        \item Perform basic error injection if requested (for health-manager demos).
        \item Remain as simple and deterministic as possible.
    \end{itemize}
    \item[Non-responsibilities] Real device drivers, interrupt simulation, or DMA.
\end{description}

\subsection{Instruments}
\begin{description}[style=nextline, leftmargin=1.5em]
    \item[Primary role] Simulated sensors and actuators.
    \item[Responsibilities]
    \begin{itemize}
        \item Produce plausible sensor data (IMU, simple GPS-like position, temperature, etc.) or accept actuator commands.
        \item Support basic failure modes (stuck value, noise, complete failure) for demonstration purposes.
        \item Be driven exclusively through the Drivers Layer.
    \end{itemize}
    \item[Non-responsibilities] High-fidelity physical models or visual simulation.
\end{description}

\subsection{Communication Interface}
\begin{description}[style=nextline, leftmargin=1.5em]
    \item[Primary role] Boundary between the flight software and the external world (``ground'').
    \item[Responsibilities]
    \begin{itemize}
        \item Accept telecommands from an external source (TCP socket, named pipe, or stdin for early demos).
        \item Deliver telemetry packets to the external world.
        \item Perform minimal framing / de-framing if needed.
        \item Remain replaceable so that different transport mechanisms can be swapped later.
    \end{itemize}
    \item[Non-responsibilities] Full CCSDS TM/TC stacks, encryption, or authentication.
\end{description}

%----------------------------------------------------------------------
\section{Message Catalogue}
%----------------------------------------------------------------------

This section defines the internal message format used on the Software Bus. The format is deliberately simpler than a full CCSDS Space Packet while still retaining the concepts that are educationally valuable for a flight-software simulation.

\subsection{Design rationale}

Internal bus messages only travel between components inside the simulation. They do not cross a radio link. Consequently they prioritise clarity, ease of construction, and low overhead over full protocol compliance.

A richer, CCSDS-inspired layout may later be introduced at the Communication Interface when messages leave the system toward the ``ground''. That framing is considered out of scope for the first version of the internal bus.

\subsection{Common message header}

Every message carried by the Software Bus contains the following fields:

\begin{center}
\begin{tabular}{@{}llp{8cm}@{}}
\toprule
\textbf{Field} & \textbf{Purpose} & \textbf{Notes} \\
\midrule
Type            & Distinguishes the nature of the message & Examples: Command, Telemetry, Health, \ldots \\
Message ID           & Routing key used by the Software Bus   & The primary value the bus matches against subscriptions \\
App ID      & Identifier of the publishing component & Useful for debugging, health monitoring and tracing \\
Timestamp \& Sequence number & Per-source (or per-topic) counter      & Helps detect gaps or reordering during demos \\
Length          & Size of the payload in bytes           & Required for variable-length payloads \\
Command field   & Action specifier for command messages  & Present when Type = Command; indicates the requested action \\
Payload         & Message-specific data                  & Interpretation depends on Topic and Type \\
\bottomrule
\end{tabular}
\end{center}

\subsection{Command field}

The Command field is kept as a distinct member of the common header rather than being buried inside the payload. Its purpose is to name the action the receiving component should perform (for example ``set mode'', ``fire thruster'', ``reset counters'', etc.).

This separation makes command handling clearer: a component can inspect the Command field first and then interpret the payload according to that action. Telemetry and health messages leave this field unused or set to a neutral value.

\subsection{What is intentionally omitted}

The following CCSDS-style elements are \emph{not} part of the internal message format:

\begin{itemize}
    \item Version number and secondary-header flags
    \item Full Application Process ID (APID) bit-packing
    \item Long telemetry secondary headers
    \item Checksums or authentication fields
\end{itemize}

These may appear later in an external packet format produced by the Communication Interface, but they are not required for component-to-component communication inside the simulation.

\subsection{Future extension points}

The header is expected to remain stable for the first complete version of the simulation. Possible later additions (explicitly optional) include a timestamp field or a simple flags word. Any such addition will be documented here before it is introduced in code.

%----------------------------------------------------------------------
\section{Scheduling \& Execution Model}
%----------------------------------------------------------------------

\textit{(To be filled once the Scheduler design is finalised.)}

%----------------------------------------------------------------------
\section{Build \& Run Instructions}
%----------------------------------------------------------------------

\textit{(To be completed when the repository skeleton exists.)}

%----------------------------------------------------------------------
\section{Demonstration Scenarios}
%----------------------------------------------------------------------

Planned scenarios (subject to refinement):

\begin{enumerate}
    \item Nominal command $\rightarrow$ actuation $\rightarrow$ telemetry path.
    \item Simple GNC closed-loop behaviour.
    \item Heartbeat loss detected by Health Manager.
    \item Memory allocation pressure / failure reporting.
    \item Instrument failure injection and resulting health alert.
\end{enumerate}

%----------------------------------------------------------------------
\section{Limitations \& Future Work}
%----------------------------------------------------------------------

\begin{itemize}
    \item Single-process, non-real-time simulation on Linux.
    \item Simplified GNC and instrument models.
    \item No redundancy, voting, or formal FDIR.
    \item No network distribution of the Software Bus.
\end{itemize}

Possible future extensions (explicitly out of current scope): multi-process version, basic RTOS integration, more realistic dynamics, CCSDS-inspired framing, etc.

%----------------------------------------------------------------------
\section{Conclusion}
%----------------------------------------------------------------------

\textit{(To be written at the end of the project.)}

\end{document}